% !TEX program = pdflatex
% Example: ACM conference paper template
% Compile with: pdflatex main.tex && bibtex main && pdflatex main.tex && pdflatex main.tex
% Requires: acmart class (tlmgr install acmart)

\documentclass[sigconf,review]{acmart}

\usepackage{graphicx}
\usepackage{amsmath,amssymb}
\usepackage{booktabs}

\begin{document}

\title{Interactive Visualization of High-Dimensional Data
       Using t-SNE and UMAP}

\author{Author Name}
\email{author@university.edu}
\affiliation{%
  \institution{University Name}
  \city{City}
  \country{Country}
}

\author{Co-Author Name}
\email{coauthor@university.edu}
\affiliation{%
  \institution{University Name}
  \city{City}
  \country{Country}
}

\begin{abstract}
We present an interactive visualization tool for
high-dimensional data that combines t-SNE and UMAP
with real-time exploration capabilities.
\end{abstract}

\ccsdesc[500]{Software and its engineering}
\ccsdesc[300]{Human computer interaction}

\keywords{dimensionality reduction, visualization, interactive systems}

\maketitle

\section{Introduction}
High-dimensional data visualization is essential for
understanding complex datasets in machine learning and
data science.

\section{Background}
t-SNE~\cite{maaten2008} and UMAP~\cite{mcinnes2018} are
two popular dimensionality reduction techniques.

\section{System Design}
Our system provides:
\begin{itemize}
\item Real-time projection updates.
\item Interactive point selection and labeling.
\item Comparative view of multiple projections.
\end{itemize}

\section{Evaluation}
User studies show that our tool reduces exploration
time by 40\% compared to static visualizations.

\section{Conclusion}
We demonstrated that interactive dimensionality
reduction visualization significantly improves
data exploration efficiency.

\bibliographystyle{ACM-Reference-Format}
\bibliography{references}

\end{document}
